\chapter*{Conclusion générale}
\addcontentsline{toc}{chapter}{Conclusion générale}
\markboth{Conclusion générale}{}

Ce travail est parti d'un constat simple~: dans le red teaming des applications LLM, on attaque avec ingéniosité mais on \emph{mesure} avec légèreté. La première version du projet en était l'illustration — un outil qui se validait sur ses propres cibles, sous le verdict d'un juge jamais éprouvé. Le recentrage qui a suivi a fait de la \emph{mesure} l'objet du mémoire.

La contribution tient en une idée et sa mise à l'épreuve. L'idée~: découpler un \textbf{oracle déterministe}, vérité-terrain objective du succès, d'un \textbf{juge LLM} dont la fiabilité est, pour une fois, \emph{mesurée} — contre l'oracle, puis contre un jugement humain établi en aveugle. La mise à l'épreuve~: une progression par paliers qui a brisé la circularité (deux cibles tierces, dont une au flag confirmé par un service externe), quantifié la fiabilité des juges (quatre sur cinq quasi indiscernables d'un humain~; un modèle \emph{safety} disqualifié), révélé que l'oracle lui-même a un angle mort fonction de la nature du critère, et démontré, par une comparaison \emph{pré-enregistrée} à PyRIT, que ce même oracle peut évaluer n'importe quel outil dans des conditions identiques.

\paragraph{Limites.} Elles sont assumées tout au long du mémoire. \textbf{Les échantillons sont petits} (19 et 18 tuples positifs)~: les intervalles de confiance, larges, interdisent les palmarès — entre les quatre bons juges, l'écart se réduit à un seul tuple, ce qui les rend interchangeables. \textbf{Le \emph{gold standard} repose sur un annotateur unique}, qui est aussi l'auteur de la méthode~: le dispositif anti-ancrage (oracle et verdicts machine masqués pendant l'adjudication) neutralise le biais d'\emph{ancrage}, mais non la conception personnelle de « ce qui compte comme vulnérabilité ». Les critères ont été \emph{pré-spécifiés} pour en borner la subjectivité~; seule une \textbf{double adjudication} indépendante, avec mesure de l'accord inter-annotateurs sur un sous-échantillon, lèverait franchement ce doute — c'est le premier ajout à prévoir. \textbf{Les deux bancs sont par ailleurs des jeux conçus pour être résolus} (un CTF et un défi public), non des applications de production~: c'est précisément ce qu'il fallait pour briser la circularité (cible tierce, oracle externe), mais cela borne la validité externe. Il faut donc distinguer nettement ce qui se transfère — \emph{la méthodologie} de découplage oracle/juge, applicable à n'importe quelle cible — de ce qui n'en sort pas — \emph{les chiffres absolus}, attachés à ces cibles CTFs. \textbf{La comparaison à PyRIT}, enfin, ne couvre qu'un orchestrateur, un attaquant et une classe de tâches qui avantage les graines directes~; elle ne dit rien des jailbreaks conversationnels, où PyRIT est dans son élément. Le \emph{harnais}, lui, reste reproductible quand modèles et cibles évoluent.

% >>> À RÉÉCRIRE TOI-MÊME (priorité HAUTE) :
% Ce paragraphe doit être à 100 % ta voix : ce que CE projet t'a appris,
% personnellement (rigueur, honnêteté des résultats, gestion d'imprévus).
\paragraph{Apport personnel.} Ce projet m'a appris à \emph{remettre en question mes propres résultats}~: chercher les angles morts , refuser de masquer un pivot imposé par une contrainte de quota, vérifier qu'un score nul n'était pas un bug. La leçon que j'en retire est la suivante : un résultat n'a de valeur que si l'on a sérieusement tenté de le réfuter.

\paragraph{Perspectives.} Plusieurs prolongements se dessinent~: une seconde adjudication humaine pour consolider la vérité-terrain~; l'élargissement de la comparaison à d'autres outils (un scanner à \emph{probes} comme garak) et à d'autres paradigmes (jailbreaks multi-tours), toujours sous le même oracle~; et une étude multi-modèles à plus grand $n$ pour resserrer les intervalles. Tous s'inscrivent dans la même logique~: non pas attaquer plus fort, mais mesurer plus juste.
